\section{总结}

\subsection{常见问题}

\begin{frame}{常见问题}
  \begin{itemize}
    \item \alert{编译不通过} 缺少必要宏包，命令拼写错误，括号未配对等
    \item \alert{表格图片乱跑} \LaTeX{} 自身的浮动定位算法
    \item \alert{段落间距变大} \LaTeX{} 排版算法
    \item \alert{参考文献} 推荐使用 \BibTeX{}，也可以手写 \cmd{bibitem}
  \end{itemize}
\end{frame}

\subsection{学习资源}

\begin{frame}{系统学习}
  \begin{itemize}
      \item 包太雷 《\LaTeX{} Notes(第二版)》~(3小时)
      \item Stefan Kottwitz 《LaTeX Cookbook》
      \item WikiBooks
        \begin{itemize}
          \item \url{https://en.wikibooks.org/wiki/LaTeX}
          \item \url{https://zh.wikibooks.org/wiki/LaTeX}
        \end{itemize}
      \item 经典文档
        \begin{itemize}
          \item 仔细阅读《一份不太简短的~\LaTeXe{} 介绍》(lshort-zh)~(1--2 天)
          \item 粗略阅读《\LaTeXe{} 插图指南》~(2--3 小时)
        \end{itemize}
      %\item 仔细阅读《\ThuThesis{} 用户手册》~(20 分钟)
      %\item 从~\ThuThesis{} 示例文档入手
  \end{itemize}
\end{frame}

\begin{frame}{扩展阅读}
  \begin{itemize}
    \item 一份其实很短的 \LaTeX 入门文档 (Liam Huang) \\
      \url{http://liam0205.me/2014/09/08/latex-introduction/} 
    \item 网站推荐: 
      \begin{itemize}
        \item http://www.latexstudio.net/
        \item http://www.chinatex.org/
      \end{itemize}
    \item 知乎专栏: http://zhuanlan.zhihu.com/LaTeX
    %\item \ThuThesis{}使用向导 v3.0 （薛瑞尼）
    \item \LaTeX{}杂谈（刘海洋）
    \item 《\LaTeX{}入门》（刘海洋）
    \item \LaTeX\ Tips:
      \url{https://alick.fedorapeople.org/fudcon-apac-2014/latex-tips.pdf}
    \item Linux 用户：\url{https://github.com/alick/fad-texlive-talk}
  \end{itemize}
\end{frame}


\begin{frame}{利用文档}
  \begin{itemize}
    \item 常用文档
      \begin{itemize}
        \item symbols: 符号大全
        \item Mathmode: 数学参考
        \item ctex, xeCJK: 中文支持
        \item texlive-zh: \TL 安装与使用
        \item 所用宏包文档
      \end{itemize}
    \item 工具
      \begin{itemize}
        \item tlmgr: \TL 管理器
        \item texdoc: \TeX{} 文档查看器\\
          例如：\texttt{texdoc lshort-zh}
        \item \url{http://texdoc.net/}
        \item TeX Studio 和 WinEdt 都支持在帮助里看文档
      \end{itemize}
  \end{itemize}
\end{frame}

\begin{frame}{一点人生的经验}
  \begin{itemize}
    \item 不要过于相信网上的中文文档
      \begin{itemize}
        \item 简单鉴别方法: 排版的好看程度
      \end{itemize}
    %\item 湿兄用U盘拷给你的的ctex套装一定是过时的，ThuThesis 八成是老版本的
    \item 如果你要处理中文
      \begin{itemize}
        \item 使用 XeLaTeX, 使用 XeLaTeX, 使用 XeLaTeX
        \item 忘记 CJK, 忘记 CJK, 忘记 CJK
        \item 使用 xeCJK
        \item 使用 ctex 宏包(2.0以上版本) (跟 C\TeX 套装仅仅是名字像)
      \end{itemize}
    \item 写一点，编译一次，减小排错搜索空间
  \end{itemize} 
\end{frame}

% 寻求帮助
\begin{frame}{求助}
  \begin{columns}[c]
    \begin{column}{.45\textwidth}
      \begin{itemize}
        \item BBS
          \begin{itemize}
            %\item \href{http://www.newsmth.net/nForum/board/TeX}{水木
             % 社区 TeX 版}
            \item \href{http://bbs.ctex.org/}{bbs.ctex.org}
          \end{itemize}
        \item \href{http://www.tex.ac.uk/cgi-bin/texfaq2html}{UK FAQ}
        \item TeX StackExchange
        \item \href{http://justfuckinggoogleit.com/}{Google}
          \begin{itemize}
            \item 使用英语搜索
          \end{itemize}
      \end{itemize}
    \end{column}
    \begin{column}{.45\textwidth}
      \includegraphics[width=\textwidth]{TFZsuperellipse-crop.pdf}
    \end{column}
  \end{columns}
\end{frame}

\begin{frame}{你也可以帮助}
  \begin{itemize}
    \item 错误反馈：GitHub Issues
    \item 改进建议：GitHub Issues
    \item 出力维护：LaTeX 宏包编写、Git
    \item 科普、答疑%\pause \hspace{2em}\alert{图书馆讲座征主讲人！}
  \end{itemize}
\end{frame}

%%% vim: set ts=2 sts=2 sw=2 isk+=\: et tw=80 cc=+1 formatoptions+=mM:
